\documentclass[12pt]{article}
\usepackage[utf8]{inputenc}
\usepackage[T1]{fontenc}
\usepackage[margin=1in]{geometry}
\usepackage{setspace}
\usepackage{xcolor}
\usepackage[colorlinks=true, urlcolor=blue!60!black, linkcolor=blue!60!black]{hyperref}
\setlength{\parindent}{2em}
\setlength{\parskip}{1.25em}
\renewcommand{\baselinestretch}{1.0}
\title{\textit{Yick Wo v. Hopkins}, 118 U.S. 356 (1886)}
\author{Case Brief}
\date{}
\begin{document}
\maketitle

\section{Case Metadata}

\begin{description}
    \item[Case Name:] \emph{Yick Wo v. Hopkins}, 118 U.S. 356 (1886) (consolidated and argued with \emph{Wo Lee v. Hopkins}, 118 U.S. 356 (1886))
    \item[Parallel Citations:] 6 S. Ct. 1064; 30 L. Ed. 220
    \item[Court:] Supreme Court of the United States
    \item[Dates:] Submitted April 14, 1886; Decided May 10, 1886
    \item[Vote:] Unanimous (9--0)
    \item[Author:] Justice Stanley Matthews (joined by Chief Justice Morrison R. Waite and Associate Justices Samuel F. Miller, Stephen J. Field, Joseph P. Bradley, John Marshall Harlan, William B. Woods, Horace Gray, and Samuel Blatchford)
    \item[Concurrences / Dissents:] None
    \item[Lower Court Dispositions:] \emph{In re Yick Wo}, 68 Cal. 294, 9 P. 139 (1885) (Supreme Court of California, discharging writ of habeas corpus and remanding petitioner); \emph{In re Wo Lee}, 26 F. 471 (C.C.D. Cal. 1886) (U.S. Circuit Court for the District of California, Circuit Judge Lorenzo Sawyer, discharging writ out of deference to the California Supreme Court despite expressing the view that the ordinance violated the Fourteenth Amendment)
    \item[Primary Opinion Sources:] \href{https://supreme.justia.com/cases/federal/us/118/356/}{Justia U.S. Supreme Court Center} $|$ \href{https://www.law.cornell.edu/supremecourt/text/118/356}{Cornell Legal Information Institute (LII)}
    \item[Briefing Book Chapter:] Equal Protection
\end{description}

\section{Facts and Procedural History}

\subsection{The San Francisco Laundry Ordinances}
In 1880, amid intense anti-Chinese political agitation in California, the Board of Supervisors of the City and County of San Francisco enacted two ordinances governing commercial laundries: Order No.~1,569 (passed May 26, 1880) and Order No.~1,587 (passed July 28, 1880). Section~1 of Order No.~1,569 and Section~68 of Order No.~1,587 contained virtually identical language making it unlawful ``for any person or persons to establish, maintain, or carry on a laundry within the corporate limits of the city and county of San Francisco without having first obtained the consent of the board of supervisors, except the same be located in a building constructed either of brick or stone.'' Section~2 of Order No.~1,569 separately barred erecting or maintaining roof scaffolding without written Board permission, and Section~3 made violations a misdemeanor punishable by a fine of up to \$1,000, imprisonment in the county jail for up to six months, or both (\href{https://www.law.cornell.edu/supremecourt/text/118/356}{118 U.S. at 357--358}).

Crucially, the ordinances established no objective fire-safety, structural, or sanitary criteria governing when the Board of Supervisors was required to grant or withhold consent. On their face, however, the ordinances made no mention of race, ancestry, or nationality.

\subsection{Stipulated Record and Stark Pattern of Enforcement}
The record before the Supreme Court contained admitted facts that revealed the practical operation of the ordinances:
\begin{itemize}
    \item \textbf{Yick Wo's Background and Compliance:} Yick Wo, a native of China who immigrated to California in 1861 and remained a subject of the Emperor of China, had operated a laundry in the same wooden building in San Francisco for twenty-two years. His premises had passed inspection by the Board of Fire Wardens (license dated March 3, 1884), which certified that his stoves, washing and drying apparatus, and iron-heating appliances were in good condition, posed no fire danger to surrounding property, and complied with San Francisco fire-safety Orders No.~1,617 and No.~1,670. The city Health Officer had likewise certified that his premises were properly drained and posed no sanitary hazard. Yick Wo held a valid city business license running through October 1, 1885. Nevertheless, when he applied to the Board of Supervisors on June 1, 1885, for consent to continue operating, the Board denied his application on July 1, 1885, without assigning any reason (\href{https://supreme.justia.com/cases/federal/us/118/356/}{118 U.S. at 358}).
    \item \textbf{Industry Demographics:} As of February 24, 1880, San Francisco had roughly 320 laundries, of which approximately 240 were owned and operated by Chinese subjects. About 310 of the 320 laundries were located in wooden buildings---the same construction material used in nine-tenths of all houses in the city. Chinese laundry operators had invested at least \$200,000 in capital and paid roughly \$180,000 annually in rent, licenses, taxes, gas, and water (\href{https://supreme.justia.com/cases/federal/us/118/356/}{118 U.S. at 358--359}).
    \item \textbf{Near-Total Racial Disparity in Permitting and Arrests:} Yick Wo and 200 other Chinese operators who had run laundries in their buildings for more than twenty years petitioned the Board for consent; \emph{every single one} of those petitions was denied, and more than 150 Chinese laundrymen were arrested for operating without consent. Conversely, of the roughly eighty non-Chinese operators running wooden-building laundries under similar conditions (qualified on the record only as lacking roof scaffolding), \emph{every petition was granted} save one---that of Mrs.~Mary Meagles---and non-Chinese operators were left unmolested to capture the displaced trade (\href{https://supreme.justia.com/cases/federal/us/118/356/}{118 U.S. at 359}).
\end{itemize}

\subsection{Procedural History}
After continuing to operate his laundry without Board consent, Yick Wo was convicted in Police Judge's Court No.~2 of San Francisco, fined \$10, and committed to the county jail by Sheriff Peter Hopkins at the rate of one day per dollar upon defaulting on the fine. In August 1885, Yick Wo petitioned the Supreme Court of California for a writ of habeas corpus. The California Supreme Court upheld the ordinances as legitimate police regulations aimed at fire safety, treating the constitutional objections as foreclosed by \href{https://supreme.justia.com/cases/federal/us/113/27/}{\emph{Barbier v. Connolly}}, 113 U.S. 27 (1885), and \href{https://supreme.justia.com/cases/federal/us/113/703/}{\emph{Soon Hing v. Crowley}}, 113 U.S. 703 (1885), and remanded Yick Wo to custody (\href{https://supreme.justia.com/cases/federal/us/118/356/}{118 U.S. at 360}).

In the companion case, \emph{Wo Lee v. Hopkins}, another Chinese laundry operator imprisoned on virtually identical facts sought federal habeas corpus relief in the U.S. Circuit Court for the District of California. Circuit Judge Lorenzo Sawyer issued a forceful opinion (\emph{In re Wo Lee}, 26 F. 471 (1886)) concluding that the ordinance covered vast sparsely populated territory (including farming lands, sand banks, and Goat Island), vested arbitrary discretion in the Board, and was plainly executed to destroy Chinese-owned laundries in violation of the Fourteenth Amendment and the 1880 treaty with China. Out of comity and deference to the California Supreme Court's ruling in \emph{Yick Wo}, however, Judge Sawyer discharged the writ and remanded Wo Lee so the U.S. Supreme Court could resolve the conflict (\href{https://supreme.justia.com/cases/federal/us/118/356/}{118 U.S. at 361--363}). Yick Wo sought review by writ of error from the California Supreme Court, and Wo Lee appealed from the federal Circuit Court; the Supreme Court heard and decided both cases together.

\section{Issue(s) Presented}

\begin{enumerate}
    \item \textbf{Scope of Fourteenth Amendment Protection:} Do the Due Process and Equal Protection Clauses of the Fourteenth Amendment protect noncitizens (specifically, resident subjects of the Emperor of China) within the territorial jurisdiction of a State?
    \item \textbf{Discriminatory Administration of a Facially Neutral Law:} Does the administration of a facially race-neutral municipal licensing ordinance violate the Equal Protection Clause of the Fourteenth Amendment when public officials exercise unguided discretion to deny permits to virtually all Chinese applicants while granting them to virtually all similarly situated non-Chinese applicants?
    \item \textbf{Unfettered Arbitrary Discretion:} Does a municipal ordinance that confers naked, standardless power upon local officials to grant or withhold permission to engage in a lawful occupation exceed constitutional bounds under the Fourteenth Amendment?
\end{enumerate}

\section{Holding}

\textbf{Yes on all issues.} Writing for a unanimous Court, Justice Matthews reversed both the Supreme Court of California in \emph{Yick Wo} and the U.S. Circuit Court in \emph{Wo Lee}, remanding each case with directions to discharge the petitioners from custody (\href{https://supreme.justia.com/cases/federal/us/118/356/}{118 U.S. at 374}). The Court held:
\begin{enumerate}
    \item The Fourteenth Amendment's guarantees of due process and equal protection apply universally to ``all persons within the territorial jurisdiction, without regard to any differences of race, of color, or of nationality,'' thereby protecting resident noncitizens equally with citizens (\href{https://supreme.justia.com/cases/federal/us/118/356/}{118 U.S. at 369}).
    \item Even where a law is ``fair on its face and impartial in appearance,'' if public authorities apply and administer it ``with an evil eye and an unequal hand'' to produce unjust discriminations between similarly situated persons based on race or nationality, that discriminatory enforcement constitutes state action violating the Equal Protection Clause (\href{https://supreme.justia.com/cases/federal/us/118/356/}{118 U.S. at 373--374}).
    \item A municipal ordinance that vests officials with naked, arbitrary power---unbounded by objective standards of applicant fitness or building safety---to forbid a lawful occupation violates the constitutional principle of a government of laws, not men (\href{https://supreme.justia.com/cases/federal/us/118/356/}{118 U.S. at 366--368, 370}).
\end{enumerate}

\section{Reasoning}

\subsection{Independent Federal Construction of the Ordinances}
Justice Matthews began by addressing federalism and appellate jurisdiction. Although the Supreme Court ordinarily defers to a state high court's construction of state law, federal constitutional review requires the Court to put its own independent construction on municipal ordinances when determining whether their enforcement violates the Federal Constitution (\href{https://supreme.justia.com/cases/federal/us/118/356/}{118 U.S. at 365--366}). Rejecting the California Supreme Court's benign reading, the Court concluded that the San Francisco ordinances did not confer ``discretion in the legal sense of that term''---exercised upon objective fire-safety conditions---but instead granted ``a naked and arbitrary power to give or withhold consent not only as to places, but as to persons.'' Even a fully qualified applicant who satisfied every safety requirement could not obtain mandamus relief, because the Board could answer that the ordinance authorized it to withhold consent ``without reason and without responsibility'' (\href{https://supreme.justia.com/cases/federal/us/118/356/}{118 U.S. at 366--367}).

\subsection{Distinguishing \emph{Barbier v. Connolly} and \emph{Soon Hing v. Crowley}}
The Court dismantled the California Supreme Court's reliance on \href{https://supreme.justia.com/cases/federal/us/113/27/}{\emph{Barbier v. Connolly}}, 113 U.S. 27 (1885), and \href{https://supreme.justia.com/cases/federal/us/113/703/}{\emph{Soon Hing v. Crowley}}, 113 U.S. 703 (1885). Those earlier San Francisco laundry decisions had upheld a uniform curfew prohibiting washing and ironing in public laundries between 10:00 p.m. and 6:00 a.m. within designated fire limits. That curfew was a genuine police regulation that treated all similarly situated operators alike under a general rule. By contrast, Orders No.~1,569 and No.~1,587 prescribed no uniform safety standard to which wooden-building operators could conform; instead, they divided wooden-building laundrymen by an arbitrary line into two classes, rendering both classes ``tenants at will, under the supervisors, of their means of living'' (\href{https://supreme.justia.com/cases/federal/us/118/356/}{118 U.S. at 367--368}). Nor did the ordinance resemble traditional tavern or liquor licensing, where statutes require an officer to make a quasi-judicial determination of personal fitness against defined legal criteria (\href{https://supreme.justia.com/cases/federal/us/118/356/}{118 U.S. at 368}).

\subsection{Universal Textual Scope of the Fourteenth Amendment (Noncitizens as ``Persons'')}
Turning to the petitioners' status as Chinese nationals, the Court emphasized two reinforcing authorities. First, Article~III of the Treaty of November 17, 1880, between the United States and China (22 Stat. 827) pledged that Chinese residents would receive the same protection, rights, and exemptions enjoyed by subjects of the most favored nation. Second, and dispositively as a matter of constitutional text, Section~1 of the Fourteenth Amendment purposefully shifts from ``citizens'' (in the Privileges or Immunities Clause) to ``any person'' (in the Due Process and Equal Protection Clauses). Reinforced by Section~1977 of the Revised Statutes (now 42 U.S.C.~\S~1981), the Equal Protection Clause is ``a pledge of the protection of equal laws'' extending universally to all persons within the territorial jurisdiction, regardless of race, color, or nationality (\href{https://supreme.justia.com/cases/federal/us/118/356/}{118 U.S. at 368--369}).

\subsection{Constitutional Antipathy to Arbitrary Personal Power}
Before turning to the enforcement record, Justice Matthews articulated a sweeping structural theory of American constitutionalism. In the American system, sovereignty resides in the people, and ``the law is the definition and limitation of power.'' Invoking the Massachusetts Bill of Rights' ideal of ``a government of laws, and not of men,'' the Court declared that compelling one person to hold life or the means of living ``at the mere will of another'' is ``the essence of slavery itself'' (\href{https://supreme.justia.com/cases/federal/us/118/356/}{118 U.S. at 369--370}). Drawing an analogy to voting rights---described in famous dicta as ``a fundamental political right, because preservative of all rights''---and citing the Maryland Court of Appeals' invalidation of a standardless steam-engine permit revocation ordinance in \emph{City of Baltimore v. Radecke}, 49 Md. 217 (1878), the Court signaled that ordinances conferring unbridled administrative discretion are suspect on their face (\href{https://supreme.justia.com/cases/federal/us/118/356/}{118 U.S. at 370--373}).

\subsection{Unconstitutional Administration: ``An Evil Eye and an Unequal Hand''}
Ultimately, the Court did not need to rely solely on facial invalidity or speculate about potential abuse, because the admitted record presented the ordinances ``in actual operation'' (\href{https://supreme.justia.com/cases/federal/us/118/356/}{118 U.S. at 373}). Even if a statute is facially neutral and impartial in text, discriminatory administration by public officials representing the State violates the Fourteenth Amendment. Here, both petitioners had satisfied every fire and health inspection, yet the Board denied consent to them and 200 other Chinese applicants while permitting roughly 80 non-Chinese operators to run wooden laundries under similar conditions. Because no legitimate non-racial explanation appeared on the record, ``the conclusion cannot be resisted that no reason for it exists except hostility to the race and nationality to which the petitioners belong'' (\href{https://supreme.justia.com/cases/federal/us/118/356/}{118 U.S. at 374}). Such administration denied equal protection and rendered the petitioners' imprisonment illegal.

\section{Concurring and Dissenting Opinions}

There were no concurring or dissenting opinions. All nine Justices joined Justice Matthews's opinion for the Court.

\section{Key Quotations}

All quotations below are verified verbatim against the primary opinion text reported at \href{https://supreme.justia.com/cases/federal/us/118/356/}{Justia (118 U.S. 356)} and \href{https://www.law.cornell.edu/supremecourt/text/118/356}{Cornell LII (118 U.S. 356)}.

\begin{enumerate}
    \item \textbf{Universal Application to Noncitizens and ``Equal Laws''} (\href{https://supreme.justia.com/cases/federal/us/118/356/}{118 U.S. at 369}):
    \begin{quote}
        ``The Fourteenth Amendment to the Constitution is not confined to the protection of citizens. It says: `Nor shall any State deprive any person of life, liberty, or property without due process of law; nor deny to any person within its jurisdiction the equal protection of the laws.' These provisions are universal in their application to all persons within the territorial jurisdiction, without regard to any differences of race, of color, or of nationality, and the equal protection of the laws is a pledge of the protection of equal laws.''
    \end{quote}

    \item \textbf{Rejection of Arbitrary Power Over Livelihood} (\href{https://supreme.justia.com/cases/federal/us/118/356/}{118 U.S. at 370}):
    \begin{quote}
        ``For the very idea that one man may be compelled to hold his life, or the means of living, or any material right essential to the enjoyment of life at the mere will of another seems to be intolerable in any country where freedom prevails, as being the essence of slavery itself.''
    \end{quote}

    \item \textbf{Voting as ``Preservative of All Rights''} (\href{https://supreme.justia.com/cases/federal/us/118/356/}{118 U.S. at 370}):
    \begin{quote}
        ``Though not regarded strictly as a natural right, but as a privilege merely conceded by society according to its will under certain conditions, nevertheless it is regarded as a fundamental political right, because preservative of all rights.''
    \end{quote}

    \item \textbf{The Canonical ``Evil Eye and Unequal Hand'' Rule} (\href{https://supreme.justia.com/cases/federal/us/118/356/}{118 U.S. at 373--374}):
    \begin{quote}
        ``Though the law itself be fair on its face and impartial in appearance, yet, if it is applied and administered by public authority with an evil eye and an unequal hand, so as practically to make unjust and illegal discriminations between persons in similar circumstances, material to their rights, the denial of equal justice is still within the prohibition of the Constitution.''
    \end{quote}

    \item \textbf{Inference of Unjustified Racial Hostility} (\href{https://supreme.justia.com/cases/federal/us/118/356/}{118 U.S. at 374}):
    \begin{quote}
        ``The fact of this discrimination is admitted. No reason for it is shown, and the conclusion cannot be resisted that no reason for it exists except hostility to the race and nationality to which the petitioners belong, and which, in the eye of the law, is not justified.''
    \end{quote}
\end{enumerate}

\section{Doctrinal Significance and Modern Legacy (Through 2026)}

\emph{Yick Wo v. Hopkins} stands as one of the foundational pillars of modern Equal Protection doctrine, though its subsequent trajectory reveals how sharply the Supreme Court has policed the boundary between \emph{discriminatory purpose} and mere \emph{disparate impact}.

\subsection{The Discriminatory-Purpose Requirement: \emph{Washington v. Davis}, \emph{Arlington Heights}, and \emph{Feeney}}
During the twentieth century, the Court invoked \emph{Yick Wo} alongside \href{https://supreme.justia.com/cases/federal/us/364/339/}{\emph{Gomillion v. Lightfoot}}, 364 U.S. 339 (1960) (invalidating an uncouth 28-sided municipal boundary redrawing that excluded almost all Black voters from Tuskegee, Alabama), as the archetype of a facially neutral law whose operation is ``unexplainable on grounds other than race.'' In \href{https://supreme.justia.com/cases/federal/us/426/229/}{\emph{Washington v. Davis}}, 426 U.S. 229 (1976), however, the Court squarely rejected a constitutional disparate-impact theory: holding that a facially neutral police entrance exam (Test 21) did not violate equal protection merely because Black applicants failed at four times the rate of white applicants, the Court ruled that challengers must prove a \emph{racially discriminatory purpose}. \emph{Davis} reframed \emph{Yick Wo} not as a disparate-impact precedent, but as an extreme case in which the statistical disparity was so stark and unexplained that discriminatory intent was inescapable.

The following Term, in \href{https://supreme.justia.com/cases/federal/us/429/252/}{\emph{Village of Arlington Heights v. Metropolitan Housing Development Corp.}}, 429 U.S. 252 (1977), the Court codified the multi-factor evidentiary framework for proving discriminatory intent behind facially neutral official action: (1) stark disparate impact (noting that cases as clear as \emph{Yick Wo} and \emph{Gomillion} are ``rare''); (2) historical background; (3) the specific sequence of events leading to the decision; (4) procedural or substantive departures from normal practice; and (5) contemporary legislative or administrative statements. Two years later, \href{https://supreme.justia.com/cases/federal/us/442/256/}{\emph{Personnel Administrator of Massachusetts v. Feeney}}, 442 U.S. 256 (1979), tightened the intent requirement further, holding that ``discriminatory purpose'' requires showing that decisionmakers acted ``because of,'' not merely ``in spite of,'' an adverse effect on an identifiable group.

\subsection{Selective Prosecution and Enforcement: From \emph{Ah Sin} to \emph{Armstrong} and \emph{Gonzalez v. Trevino}}
Because \emph{Yick Wo} invalidated discriminatory administration of a criminal ordinance, it became the doctrinal fountainhead for selective-prosecution and selective-enforcement claims. Yet the Court quickly erected steep evidentiary hurdles. In \href{https://supreme.justia.com/cases/federal/us/198/500/}{\emph{Ah Sin v. Wittman}}, 198 U.S. 500 (1905), the Court rejected a Chinese defendant's \emph{Yick Wo} challenge to a San Francisco anti-gambling ordinance because he failed to prove that non-Chinese violators existed and were not prosecuted. Nearly a century later, in \href{https://supreme.justia.com/cases/federal/us/517/456/}{\emph{United States v. Armstrong}}, 517 U.S. 456 (1996), the Court cited \emph{Yick Wo} while imposing a demanding threshold even to obtain \emph{discovery} on a selective-prosecution claim: a criminal defendant must produce ``some evidence'' of both discriminatory effect (that similarly situated individuals of a different race were \emph{not} prosecuted) and discriminatory intent, overcoming a strong presumption of prosecutorial regularity.

In the First Amendment retaliatory-arrest context, the Court borrowed directly from \emph{Armstrong} and \emph{Yick Wo}'s comparator logic in \href{https://supreme.justia.com/cases/federal/us/587/17-1174/}{\emph{Nieves v. Bartlett}}, 587 U.S. 391 (2019), creating an exception to the no-probable-cause bar when a plaintiff presents objective evidence that similarly situated persons not engaged in protected speech were not arrested. In \href{https://supreme.justia.com/cases/federal/us/602/22-1025/}{\emph{Gonzalez v. Trevino}}, 602 U.S. 653 (2024), the Court clarified that a plaintiff need not produce ``virtually identical'' comparator arrests to satisfy \emph{Nieves}; statistical evidence showing a statute had never been used to arrest anyone for similar conduct over a decade could suffice---an approach echoing \emph{Yick Wo}'s willingness to infer improper administration from an anomalous enforcement record.

\subsection{Equal Protection Rights of Noncitizens}
\emph{Yick Wo}'s holding that ``any person'' includes noncitizens launched an enduring line of alienage jurisprudence, including \href{https://supreme.justia.com/cases/federal/us/239/33/}{\emph{Truax v. Raich}}, 239 U.S. 33 (1915) (striking down an Arizona quota restricting alien employment), \href{https://supreme.justia.com/cases/federal/us/403/365/}{\emph{Graham v. Richardson}}, 403 U.S. 365 (1971) (applying strict scrutiny to state alienage classifications denying welfare benefits), and \href{https://supreme.justia.com/cases/federal/us/457/202/}{\emph{Plyler v. Doe}}, 457 U.S. 202 (1982) (confirming that even undocumented children are ``persons within [a State's] jurisdiction'' protected against denial of public K--12 education). At the same time, the Court carved out major exceptions: federal alienage classifications receive deferential review under Congress's plenary immigration power (\href{https://supreme.justia.com/cases/federal/us/426/67/}{\emph{Mathews v. Diaz}}, 426 U.S. 67 (1976)), and states may restrict certain ``political function'' offices (such as police officers and public school teachers) to citizens (\href{https://supreme.justia.com/cases/federal/us/441/68/}{\emph{Ambach v. Norwick}}, 441 U.S. 68 (1979)).

\subsection{Contemporary Equal Protection Doctrine Through 2026}
Through 2026, \emph{Yick Wo} continues to anchor both anti-classification rhetoric and high-stakes litigation over facially neutral policies:
\begin{itemize}
    \item \textbf{Universal Anti-Discrimination Principle:} In \href{https://supreme.justia.com/cases/federal/us/600/20-1199/}{\emph{Students for Fair Admissions, Inc. v. President and Fellows of Harvard College}}, 600 U.S. 181 (2023), Chief Justice Roberts cited \emph{Yick Wo} early in the opinion to establish that the Fourteenth Amendment's ``broad and benign provisions'' apply universally ``to all persons'' and condemn all unjustified ``hostility to . . . race and nationality''---framing the Clause's core command as colorblind equality.
    \item \textbf{Proof of Intent in Executive, Redistricting, and Legislative Challenges:} Modern challengers invoking \emph{Yick Wo} and \emph{Arlington Heights} face formidable evidentiary burdens. In \href{https://supreme.justia.com/cases/federal/us/591/18-587/}{\emph{Department of Homeland Security v. Regents of the University of California}}, 591 U.S. 1 (2020), a plurality rejected an Equal Protection challenge to the rescission of DACA despite its overwhelming impact on Latino recipients, finding the allegations insufficient under \emph{Arlington Heights}. Similarly, in \href{https://supreme.justia.com/cases/federal/us/602/22-807/}{\emph{Alexander v. South Carolina State Conference of the NAACP}}, 602 U.S. 1 (2024), the Court reinforced a strong ``presumption of legislative good faith'' in racial-gerrymandering challenges, requiring plaintiffs to disentangle race from partisan politics and faulting them for not producing an alternative map. And in \href{https://supreme.justia.com/cases/federal/us/605/23-477/}{\emph{United States v. Skrmetti}}, 605 U.S. \_\_\_ (2025), the Court held that Tennessee's restriction on certain medical treatments for minors classified on the basis of age and medical use rather than sex or transgender status, underscoring how decisively the threshold classification-versus-intent inquiry controls modern Equal Protection outcomes.
\end{itemize}

\section{Pedagogical Critiques and Exam Hypotheticals}

\subsection{Pedagogical Critiques and Doctrinal Tensions}
\begin{itemize}
    \item \textbf{Facial Arbitrariness vs.\ As-Applied Animus:} Justice Matthews blended two distinct theories: a structural objection to ``naked and arbitrary power'' over a lawful occupation and an Equal Protection objection to racially discriminatory enforcement. Later decisions largely abandoned the first strand for ordinary economic regulation (compare \href{https://supreme.justia.com/cases/federal/us/330/552/}{\emph{Kotch v.\ Board of River Port Pilot Commissioners}}, 330 U.S. 552 (1947)) while canonizing the second for suspect classifications.
    \item \textbf{An Evidentiary Unicorn:} Because San Francisco stipulated to its 200-to-0 versus 80-to-1 permitting record and offered no safety rationale, \emph{Yick Wo} is an outlier. By treating \emph{Yick Wo} as the benchmark for inferring discriminatory purpose from disparate impact (\href{https://supreme.justia.com/cases/federal/us/426/229/}{\emph{Washington v.\ Davis}}, 426 U.S. 229 (1976); \href{https://supreme.justia.com/cases/federal/us/429/252/}{\emph{Arlington Heights}}, 429 U.S. 252 (1977); \href{https://supreme.justia.com/cases/federal/us/517/456/}{\emph{United States v.\ Armstrong}}, 517 U.S. 456 (1996)), modern doctrine makes subtler or mixed-motive discrimination far harder to prove.
    \item \textbf{Exclusion-Era Irony:} Decided four years after the Chinese Exclusion Act of 1882, \emph{Yick Wo} robustly protected resident noncitizens' livelihoods against local expropriation even as the Court deferred broadly to federal exclusion and deportation powers.
\end{itemize}

\subsection{Exam Hypothetical 1: Discretionary Permits and Stark Disparity}
\textbf{Fact Pattern:} The City of Pacifica requires all street-food commissaries in pre-1970 buildings to obtain a ``Fire Resilience Permit'' from a Licensing Board empowered to grant or deny applications ``consistent with neighborhood safety,'' with no objective standards. Ninety percent of pre-1970 commissaries are operated by Central American immigrants (many noncitizens). Over two years, the Board denies 142 of 145 applications from Central American operators---including 110 with clean fire-marshal inspections---while granting 38 of 40 applications from non-Hispanic operators in comparable buildings, citing only ``neighborhood congestion.''

\textbf{Analysis:} Under \href{https://supreme.justia.com/cases/federal/us/118/356/}{\emph{Yick Wo}} (118 U.S. at 369, 373--374), noncitizens are ``persons'' protected by the Equal Protection Clause. The extreme disparity (97.9\% vs.\ 5\% denial rates), clean safety inspections, and standardless discretion establish administration with ``an evil eye and an unequal hand,'' supporting an inference of purposeful national-origin or alienage discrimination under \emph{Davis} and \emph{Arlington Heights} that fails strict scrutiny.

\subsection{Exam Hypothetical 2 (``Opposite Side'' Contrast): Uniform Objective Safety Code}
\textbf{Fact Pattern:} Pacifica instead repeals the permit system and requires \emph{every} commercial commissary citywide to install a UL-listed \$45,000 fire-suppression system, with no waivers and purely mechanical inspection. Because immigrant operators disproportionately lease smaller, older buildings, 85\% of Central American-owned commissaries close versus 12\% of corporate-backed commissaries. Council minutes show actuaries recommended the exact specification after three fatal kitchen fires, with no evidence of racial or anti-immigrant animus.

\textbf{Analysis:} This regulation tracks \href{https://supreme.justia.com/cases/federal/us/113/27/}{\emph{Barbier v.\ Connolly}}, 113 U.S. 27 (1885), and \href{https://supreme.justia.com/cases/federal/us/113/703/}{\emph{Soon Hing v.\ Crowley}}, 113 U.S. 703 (1885), rather than \emph{Yick Wo}. Under \emph{Davis} and \href{https://supreme.justia.com/cases/federal/us/442/256/}{\emph{Personnel Administrator of Massachusetts v.\ Feeney}}, 442 U.S. 256, 279 (1979), foreseeable disparate impact alone does not trigger strict scrutiny unless Pacifica acted ``because of,'' not merely ``in spite of,'' its effect on immigrants. Supported by an objective fire-safety record and uniform enforcement, the code survives rational basis review.

\end{document}
